% Figura: integrador inversor (figura 5a del guion) y derivador inversor (figura 5b).
\begin{figure}[ht]
\centering
\begin{minipage}[b]{0.48\textwidth}\centering
\begin{circuitikz}[circuito, scale=0.85, transform shape]
  \draw (0,0) node[op amp, ao] (A) {};
  \draw (A.-) -- ++(-0.5,0) coordinate (N);
  \draw (N) to[R, l_=$R$, i<^=$i$] ++(-2.2,0) to[short, -o] ++(-0.3,0) node[left, tens] {$\vin$};
  \draw (A.out) -- ++(0.6,0) coordinate (O);
  \draw (N) -- ++(0,1.2) coordinate (T1) to[C, l_=$C$] (T1 -| O) -- (O);
  \draw (T1) -- ++(0,1.0) coordinate (T2) to[R, l=$R_F$] (T2 -| O) -- (T1 -| O);
  \fill (T1) circle (1.6pt); \fill (T1 -| O) circle (1.6pt);
  \draw (O) to[short, *-o] ++(0.6,0) node[right, tens] {$\vout$};
  \draw (A.+) -- ++(-0.5,0) -- ++(0,-0.4) node[ground] {};
  \fill (N) circle (1.6pt);
\end{circuitikz}\\[2pt]
{\small (a)}
\end{minipage}\hfill
\begin{minipage}[b]{0.48\textwidth}\centering
\begin{circuitikz}[circuito, scale=0.85, transform shape]
  \draw (0,0) node[op amp, ao] (A) {};
  \draw (A.-) -- ++(-0.5,0) coordinate (N);
  \draw (N) to[C, l_=$C$, i<^=$i$] ++(-2.2,0) to[short, -o] ++(-0.3,0) node[left, tens] {$\vin$};
  \draw (A.out) -- ++(0.6,0) coordinate (O);
  \draw (N) -- ++(0,1.2) coordinate (T1) to[R, l=$R$] (T1 -| O) -- (O);
  \draw (O) to[short, *-o] ++(0.6,0) node[right, tens] {$\vout$};
  \draw (A.+) -- ++(-0.5,0) -- ++(0,-0.4) node[ground] {};
  \fill (N) circle (1.6pt);
\end{circuitikz}\\[2pt]
{\small (b)}
\end{minipage}
\caption{(a) Integrador inversor (figura~5a del guion); con $R_F\to\infty$ (set~2) es un integrador ideal.
(b) Derivador inversor (figura~5b del guion).}
\label{fig:int-der}
\end{figure}
